\section{The common algebraic family for the residue comparison}\label{sec:deformations}

Let $C$ be a smooth projective curve of genus at least two and $B$ a
smooth projective variety. We start with
\[
 X_0=\operatorname{Tot}\big((L\boxtimes P)\oplus(M\boxtimes Q)\big)
       \longrightarrow C\times B,
 \qquad L\otimes M=K_C ,
\]
where $P,Q$ are line bundles on $B$. Fix a degree $m>0$ over $C$.
By the support of a cycle we mean the union of its distinct image
curves. If $A=\sum_i e_i A_i$ is a cycle, its reduced cycle is
$\sum_i A_i$. A \emph{proper multiple} is a cycle $eA'$ with
$e>1$ and $A'$ effective. Necessarily $e\mid m$.

Let $\tau$ be a tuple of the coefficient evaluations of
\S\ref{sec:weights}, translated by a section of its test bundle
$T\to B$. For this section it suffices that the evaluations be
linear combinations of symmetric functions of the intersection with
finitely many framed fibers over $F\subset C$, that the translation
be smooth, and that $\operatorname{rk}_{\mathbb C}T=\dim_{\mathbb C}B$.
We allow the supported Thom class to be localized in an open set of the
degree-$m$ cycle space, provided its zero set there is compact. All
motions below refer to that localized class. They preserve any scalar
factor in its normalization.

\begin{theorem}\label{thm:separate}
Suppose the compact zero set of $\tau$ contains no proper multiple.
For every $K$ there is a smooth quasi-projective variety $S$ smooth
over $B$, an algebraic section $i:B\to S$, and a smooth family of
Zariski locally trivial fibrations $X\to C\times S$ with fiber $\mathbb A^2$ restricting along $i$ to $X_0$.
There is a homotopy of the parameter section and the tests, with compact
joint zero set in the chosen cycle neighborhoods, with the following
properties.
\begin{enumerate}
\item At its endpoint, the zeros with connected support and sum of
normalization genera at most $K$ are finitely many smooth embedded
curves of multiplicity one. The relative space of embedded curves is
smooth of dimension $\dim S$ there. The target variations surject onto the normal obstruction space
there. The relative space meets the parameter graph transversely, and
the tests have transverse zero on this intersection.
\item On the single algebraic family $X/S$, the products of the test
Thom class with the normal Thom classes of the parameter graphs have
equal compactly supported classes at the two ends. At the initial end,
capping the relative stable-map virtual class with the graph class
restricts it to the relative stable-map class of $X_0/B$.
\end{enumerate}
In particular every connected stable map of arithmetic genus at most
$K$ on the endpoint zero set has one of the indicated image cycles.
\end{theorem}

The tests in this statement may have zeros of disconnected support.
The compact support in the second assertion includes such zeros. We
impose connectedness only in determining the image of a connected
stable map. We first describe the target deformations used in the
proof.

\subsection{Deformations at finitely many fibers}

Consider a family already constructed over a smooth quasi-projective
$S\to B$. Off a finite subset $P_0\subset C$, disjoint from $F$,
keep its identification with the pullback of $X_0$. Initially $S=B$
and $P_0$ is empty. Fix $p\notin F\cup P_0$. On a Zariski neighborhood
$U$ of $p$ choose frames for $L,M$ and a uniformizer $t$ regular
there with its only zero at $p$. Shrink $U$ to avoid the old punctures.
In these frames the fiber variable $y$ takes values in the pullback
of the plane bundle $E=P\oplus Q$ over $B$.

Let $b,r$ be positive integers. On $U\setminus\{p\}$ use the change
of coordinates
\begin{equation}\label{eq:shear-gluing}
 y\longmapsto y+(\iota_\ell\Omega)
           \sum_{a=1}^b\sum_{v=0}^{r-1} c_{av}t^{-a}\ell(y)^v,
 \qquad
 \ell\in E^\vee,\quad \Omega\in\bigwedge^2 E,\quad c_{av}\in\mathbb C .
\end{equation}
All the displayed parameters are coordinates on a vector bundle over
$S$. Since $\ell(\iota_\ell\Omega)=0$, this automorphism preserves
$\ell(y)$, has inverse obtained by negating the $c_{av}$, and has
Jacobian determinant one in the fiber. At $p$ we glue a new copy of
the original affine chart by a composition of two such automorphisms
with independent parameters. For a finite set of new punctures use
the old target on their complement. Each new chart is chosen to miss
every other puncture. On the intersection of two new charts use their
transitions through the old chart. This gives an algebraic family
$X'\to C\times S'$ with $S'$ the total space of the new parameter
bundle. It restricts to the old family whenever all $c_{av}=0$,
including at the algebraic zero section $S\to S'$. Off the union of
the old and new punctures the original identification is preserved.
Since the inserted determinants are one, the relative tangent
determinant over $C$ in each threefold is still pulled back from
$L\otimes M$. Every compact curve in it has threefold Chern number
zero. The affine-plane fibers contain no compact curves.

Here is the infinitesimal effect. Let
$h_i:D_i\to X_s$ be the smooth normalizations of the distinct
positive components of a cycle, with finite projections
$f_i:D_i\to C$. Set $D=\coprod_iD_i$, $f=\coprod_i f_i$, and let
$V|_{D_i}=h_i^*T_{X_s/C}$. Over a general $p$ the covers are
unramified and all $r=\sum_i\deg f_i$ sheet values in the affine
plane are distinct. Indeed the normalized maps are birational onto
different curves. The number $r$ counts the sheets of the reduced
cycle. Choose two independent forms $\ell$ which each separate their
values at $p$, and use a nonzero $\Omega$. These choices exist since
only finitely many lines in $E^\vee$ are forbidden.

\begin{lemma}\label{lem:shear-jets}
Let $Z$ be a finite subscheme of the smooth domains, disjoint from
$f^{-1}(p)$. For sufficiently large $b$, variations of the
coefficients $c_{av}$ in the two factors of \eqref{eq:shear-gluing},
at $c=0$, span the obstruction space $H^1(D,V(-Z))$ for deforming
the graphs with fixed covers and prescribed restriction to $Z$.
\end{lemma}

\begin{proof}
Since $f^*p$ has positive degree on each component, we can arrange
\[
 H^1(D,V(bf^*p)(-Z))=0 .
\]
The connecting morphism from
principal parts of order $b$ over $p$ then surjects onto
$H^1(D,V(-Z))$. On each sheet, multiplication by $t^b$ identifies
these principal parts with two copies of
$R_b=\mathbb C[t]/(t^b)$. Write $y_1(t),\ldots,y_r(t)$ for the
formal sheet values. For either $\ell$ the matrix
$(\ell(y_j(t))^v)_{1\le j\le r,\ 0\le v<r}$ has Vandermonde
determinant a unit in $R_b$. For each $v$ the coefficients
$\sum_{a=1}^b c_{av}t^{b-a}$ run through $R_b$. Thus one shear
prescribes arbitrary principal parts on every sheet in the direction
$\iota_\ell\Omega$. The two directions span the plane. Differentiation
of the transition on the punctured overlap gives exactly the resulting
connecting class, up to the choice of sign for the \v Cech
differential. Prescribing the graph on $Z$ replaces $V$ by $V(-Z)$
because the difference of two local lifts vanishes on $Z$.
\end{proof}

We will need this lemma on a compact set of smooth normalization
tuples, simultaneously for $Z$ of lengths zero, one and two.
The corresponding relative symmetric powers of the disjoint smooth
source curves are proper over the tuple space. For each tuple
one can take three good punctures, so that at least one avoids the
base images of any such $Z$. At a given puncture $b$ can be chosen
uniformly over these subschemes by relative Serre vanishing for the
ideal of the universal divisor. Surjectivity of the graph deformation
problem with the jets imposed is open by the rank condition on a two-term
resolution of the sections of $V(-Z)$. Its differential augmented by
the parameter variations is surjective. This remains true where the dimensions of
its kernel or of $H^1(D,V(-Z))$ change. Compactness gives finitely
many choices of directions on neighborhoods of the tuples and their
marked subschemes. Distinct choices of $\ell,\Omega$ at the same
Borel parameter can be placed in separate factors of the parameter
bundle. Extend each choice locally over $B$ and then smoothly by a
cutoff. With the $c_{av}$ still zero, moving the parameter section to
these reference values leaves the target family unchanged. Coincident
choices of new punctures can also be accommodated by composing their
shears in the same transition. The linearized variations add at $c=0$.

\subsection{Specialization of the weighted images}

We give the compactness argument for the tuples to which the lemma
will be applied. Each $X'\to C\times S'$ is affine of finite type by
Zariski descent, and $C\times S'$ is quasi-projective. Hence $X'$ is
quasi-projective, and admits a projective closure over $S'$. In each
threefold fiber the projection to $C$ has contractible fibers. Thus
on degree-$m$ curves any chosen relative polarization has fixed
degree locally in the parameter. On compact sets of parameters these
degrees are bounded. The degree of a positive summand is bounded by
that of the whole cycle for the polarization.

Choose a relatively compact open neighborhood of the original compact
zero set, with closure in the given open locus of cycles. Such
neighborhoods can be taken in the fixed-degree relative Chow space
of a projective closure, on the open set of cycles with support in
the affine target. They can also be obtained from a compact bound
on supports and masses \cite[Theorem~2.11]{DIW}. For $e\mid m$ the
image of multiplication by $e$ is closed in this locus. Along a
convergent sequence of multiples the roots have the same supports
and their masses are divided by $e$. A subsequential limit of the
roots is therefore a root of the limiting cycle. Choose the closure
disjoint from these images for all $e>1$. Its boundary misses the
zero set. Whenever the parameter family is enlarged, take such a
neighborhood of the whole current compact zero set in the enlarged
relative cycle space, with closure on the old graph inside the
preceding neighborhood. The compact boundary misses the simultaneous
zero of the old graph and tests. It continues to miss this zero for
all sufficiently small motions of both. The supported Thom product
in the new neighborhood thus restricts to the previous product by
excision. Applied successively, this construction keeps compact
support also for zeros whose normalizations exceed the genus cutoff.

For a cycle $A=\sum_{i=1}^r e_iA_i$ write $d_i$ for the projection
degree and $h_i$ for the normalization genus of $A_i$. Set
\[
       \lambda(A)=\left(\sum_i d_i,\ \sum_i h_i,\ -r\right)
\]
with the lexicographic order. The first entry uses the reduced cycle.
Only finitely many values occur for $\sum_i h_i\le K$.

\begin{lemma}\label{lem:weighted-specialization}
For a convergent sequence in the chosen neighborhood with fixed
$\lambda$, the limiting value of $\lambda$ is at most that of the
sequence, and the sum of normalization genera cannot increase. If
$\lambda$ stays equal, then, after relabeling and taking a
subsequence, the normalization maps converge to the smooth
normalizations of the limiting cycle, with their original weights.
Connected supports have connected limit.
\end{lemma}

\begin{proof}
Keep the degrees, genera and weights of the individual labels constant
along a subsequence. Stable-map compactness applied to their
unweighted normalization maps gives a stable limit of each label.
One can take these limits in the projective closure. The bounds on the
component degrees give properness there, and the limiting total support
lies in the affine open. The weighted sum of their image cycles is
the given limit. Write $A'_j$ for its distinct image curves and
$b_{ij}$ for the sum of the covering degrees onto $A'_j$ in the
$i$th stable limit. With $d'_j=\deg(A'_j/C)$ one has
\[
       d_i=\sum_j b_{ij}d'_j,\qquad e'_j=\sum_i e_i b_{ij}.
\]
Every column sum of $(b_{ij})$ is positive. It follows that the
reduced degree drops unless every column sum is one. Each $A'_j$
has normalization genus at least two since it covers $C$. By
Riemann--Hurwitz any component covering it has at least its genus.
The genus formula for the nodal sources now bounds the sum of the
distinct normalization genera by $\sum_i h_i$.

At equal reduced degree every positive component maps with degree
one onto a different image, and each old label has at least one
positive component. Thus the number of distinct images cannot
decrease. If the entire triple is unchanged, each label has only one
positive component, with its full genus. The dual graph is a tree and
all contracted vertices are rational. Any nontrivial such tree has a
contracted leaf, contrary to stability without markings. The
limiting maps are therefore smooth and birational onto their images,
and retain the corresponding weights. Finally, a Hausdorff limit of
nonempty compact connected sets is connected.
\end{proof}

\subsection{Separation through a fixed genus}

\begin{proof}[Proof of Theorem~\ref{thm:separate}]
Consider the current constrained cycles of connected support whose
normalization genera sum to at most $K$. Suppose the values of
$\lambda$ below a given value have been treated, with finitely many
transverse multiplicity-one embedded zeros retained. Around each
of these the embedded smooth maps form an open set of the full cycle
space \cite[Proposition~2.21]{DIW}. After any one enlargement of
the family the implicit function theorem preserves this isolated zero
with its sign under sufficiently small motions of the graph and tests.
Take smaller disjoint neighborhoods of the retained zeros.

On the complement, take the least remaining value of $\lambda$. The
normalization tuples with this value over the current constrained set
form a compact set up to finite relabeling. Indeed their total cycles
have compact closure as above. Lemma~\ref{lem:weighted-specialization}
places the limit in the same bounded connected locus. Minimality
forces equality of the triples and hence convergence in the smooth
normalization charts. The same argument along sequences of
perturbations tending to the current graph has two consequences. An
earlier absent value cannot occur on this complement under sufficiently
small perturbations. Every new zero of the present value must have
its tuple in the chosen neighborhoods of the compact original tuples.
If there is no remaining value, the induction ends.

Enlarge the family using Lemma~\ref{lem:shear-jets} on this compact
set. The space of ordered smooth graph maps over the enlarged
parameter $S'$ is smooth near the tuples. To see this, first deform
their finite covers of $C$ in the smooth Hurwitz stacks. For a smooth
finite source cover $f_i$, the cokernel of
$T_{D_i}\to f_i^*T_C$ has finite length and the kernel vanishes.
After these cover deformations and the parameter in $B$ have been
prescribed, the remaining obstruction for the graphs lies in
$\bigoplus_i H^1(D_i,V_i)$, spanned by the parameter variations.
Thus the graph space submerses onto the product of the Hurwitz
stacks and $B$. It has complex dimension $\dim S'$, since the
unmarked threefold index of each component is its Chern number zero
and the target parameters occur only once. Each graph has trivial
stabilizer by birationality. This calculation only counts distinct
normalized components. Their weights enter the tests through the
cycle sum.

Mark a critical point of a cover. This marking is finite over its
Hurwitz stack, while nonimmersion requires the two vertical first
derivatives of the graph to vanish. Mark instead two distinct source
points over the same base value, on the same or on different
components. This adds at most one complex dimension over the covers.
Coincidence of the two images then imposes two vertical equations.
Take smooth strata of these marked spaces over the Hurwitz stacks
and pull back by the submersion just proved. The fixed-cover jet
variations of Lemma~\ref{lem:shear-jets} are transverse to the
equations on every such stratum at a zero. In the second calculation
a common change of target coordinates cancels in the variation of the
difference of two equal images. In the first calculation it preserves
the vanishing of the derivative at a critical point. Hence both bad
loci have dimension strictly less than $\dim S'$. The length-two
subschemes in the lemma include the limits of two marks approaching
one another.

Move the parameter graph by finitely many smooth sections spanning
its vertical normal directions along the compact base $B$. Move the
tests independently by smooth sections spanning $T$ there. The graph
condition and the tests together have real codimension $2\dim S'$,
and the auxiliary directions span their joint normal bundle even on
the marked strata. Sard's theorem on countably many smooth charts
gives arbitrarily small choices transverse on the unmarked spaces and
missing all the bad strata. We can impose transversality to the
parameter graph alone at the same time. A connected constrained tuple
of the present type now has one component, since distinct embedded
images in the tuple are disjoint. Its weight is one, since the collar
excludes proper multiples. The surviving zeros of this type are
finite by Lemma~\ref{lem:weighted-specialization}, compactness, and
transversality in the smooth charts. They are isolated in the full
cycle space on the parameter graph with the tests imposed, again by
embedded-cycle openness. Add them to the retained list. An enlargement
preserves regularity of the parameterized graph equations at the old
zeros, since the old parameter directions are still present at its
zero section. Openness then gives the same regularity at their
continuations. At an embedded curve the obstruction for the fixed
target is $H^1$ of its normal bundle. Since the parameterized graph
problem has zero obstruction, its target variations surject onto
this group. There are only finitely many values of $\lambda$ to
treat.

Take $S$ to be the final parameter scheme. Its algebraic zero sections
embed every preceding family in this last family. Concatenating the
smooth parameter paths gives the claimed homotopy. The reference
values of $\ell,\Omega$ can be moved before each small perturbation
inside $c=0$, where the target and tests are constant. The compact
joint support there is just transported along the reference path.
The successive collars and excision give a homotopy of the joint
Thom product, relative to the complement of its compact support.
A connected stable map of genus at most $K$ has connected support,
and its arithmetic genus bounds the sum of the genera of the distinct
image normalizations by Riemann--Hurwitz. It consequently maps to
one of the retained embedded curves with degree one.

Finally, maps to fibers of the smooth open $X/S$ form an open
substack of the relative stable-map stack for a projective closure.
The obstruction theory on this open, relative to the prestable source
and to $S$, is dual to the derived sections of the pullback of
$T_{X/S}$. It gives the vertical virtual class in Borel--Moore homology
after taking cycle classes. The universal domain and this obstruction
theory commute with base change at the regular algebraic section $i$.
Its refined Gysin pullback gives the relative virtual class of the
original family \cite[\S7]{BF}. The normal Thom class computes the
same Gysin map on cycle classes. For each connected source genus the inverse
image of the compact supporting set on cycles is compact by
properness of the image-cycle map \cite{WueProduct}. This proves the second
assertion and completes the proof.
\end{proof}
